\chapter{Reduced States, Entanglement Entropy, and Unitary Evolution}

A density matrix first appears as the natural description of incomplete preparation: a quantum system has been prepared, but we do not know exactly how. Entanglement gives the same mathematical object a deeper role. Even when the state of a composite system is known completely, a subsystem generally has no state vector of its own. Its reduced density matrix contains everything that can be predicted locally. We shall derive that statement in components, test it on two spin states, and use entropy to distinguish product states from entangled ones. Then we shall turn to a second question: how a quantum state changes with time. Preservation of inner products will lead to unitarity, a Hermitian generator, and the Schr\"odinger equation.

\section{Density Matrices and Incomplete Preparation}

Suppose an electron has been placed in a strong magnetic field and allowed to align with it. Someone else chose the field direction and prepared the electron, but did not tell us the direction. The spin may therefore be in a perfectly definite pure state even though we do not know which state it is. We may know a probability distribution over possible preparations, or we may know nothing at all. A density matrix \(\rho\) is the quantum description of precisely this kind of incomplete information.

The classical analogue is a probability density on phase space or a probability distribution over discrete configurations. The quantum object is an operator with three essential properties:
\begin{equation}
\rho^\dagger=\rho,
\qquad
\operatorname{Tr}\rho=1,
\qquad
\lambda_i\geq 0,
\end{equation}
where the \(\lambda_i\) are the eigenvalues of \(\rho\). Hermiticity makes the eigenvalues real and supplies an orthonormal eigenbasis. Positivity makes them admissible probabilities, and normalization gives
\begin{equation}
\sum_{i=1}^{n}\lambda_i=1.
\end{equation}
In the eigenbasis of \(\rho\), \(\lambda_i\) is the probability assigned to the \(i\)-th eigenstate.

Complete ignorance is the case in which no direction in Hilbert space is preferred:
\begin{equation}
\rho_{\mathrm{max}}=\frac{1}{n}I.
\end{equation}
Every eigenvalue is \(1/n\). At the opposite extreme, complete knowledge is represented by a normalized state \(\ket{\psi}\), with
\begin{equation}
\rho_\psi=\ket{\psi}\bra{\psi}.
\end{equation}
This is the projector onto the one-dimensional subspace spanned by \(\ket{\psi}\).

\begin{classroomqa}
\lecturetimestamp{00:06:47}
\audiencequestion{Can a preparation be described as a fifty-fifty choice between two states that are not orthogonal, for example two spin directions separated by \(45^\circ\)?}
\lecturerresponse{Such a preparation procedure is possible, but those two weights are not generally the eigenvalues of the resulting density matrix. One first forms the weighted sum of projectors and then diagonalizes it. The clean probability interpretation of the eigenvalues belongs to the orthonormal eigenbasis of \(\rho\), not necessarily to the nonorthogonal states named by the preparer.}
\end{classroomqa}

\begin{classroomqa}
\lecturetimestamp{00:07:47}
\audiencequestion{What distinguishes a pure state from a mixed state?}
\lecturerresponse{A pure-state density matrix has one nonzero eigenvalue, equal to \(1\). A mixed state has more than one nonzero eigenvalue. Older literature used the German word \emph{Gemisch}; the modern term is simply \emph{mixed state}.}
\end{classroomqa}

Let us verify the spectrum of the pure-state projector. Normalization gives
\begin{equation}
\rho_\psi\ket{\psi}
=\ket{\psi}\braket{\psi}{\psi}
=\ket{\psi},
\end{equation}
so \(\ket{\psi}\) is an eigenvector with eigenvalue \(1\). If \(\ket{\phi}\) is orthogonal to \(\ket{\psi}\), then
\begin{equation}
\rho_\psi\ket{\phi}
=\ket{\psi}\braket{\psi}{\phi}
=0.
\end{equation}
Every direction orthogonal to \(\ket{\psi}\) therefore has eigenvalue zero.

\begin{figure}[htbp]
\centering
\includegraphics[width=0.72\textwidth]{lecture_08_figure_06.png}
\caption{The pure-state density matrix acts as a projector: \(\rho\ket{\psi}=\ket{\psi}\).}
\lectureframe{00:09:32}
\end{figure}

\section{Expectation Values from the Trace}

The operational rule for a density matrix is
\begin{equation}
\overline{M}=\operatorname{Tr}(\rho M),
\end{equation}
where \(M\) is an observable. For two factors the cyclic property of the trace gives
\begin{equation}
\operatorname{Tr}(\rho M)=\operatorname{Tr}(M\rho),
\end{equation}
whether or not \(\rho\) and \(M\) commute.

In any orthonormal basis \(\{\ket{i}\}\),
\begin{equation}
\operatorname{Tr}X=\sum_i\bra{i}X\ket{i}.
\end{equation}
The numerical value does not depend on the basis. For a pure state,
\begin{align}
\overline{M}
&=\operatorname{Tr}\!\left(\ket{\psi}\bra{\psi}M\right) \\
&=\sum_i\braket{i}{\psi}\bra{\psi}M\ket{i} \\
&=\bra{\psi}M\left(\sum_i\ket{i}\bra{i}\right)\ket{\psi} \\
&=\bra{\psi}M\ket{\psi}.
\end{align}
The resolution of the identity,
\begin{equation}
\sum_i\ket{i}\bra{i}=I,
\end{equation}
turns the density-matrix rule back into the familiar pure-state expectation value.

\begin{figure}[htbp]
\centering
\includegraphics[width=0.74\textwidth]{lecture_08_figure_07.png}
\caption{The trace calculation for a pure-state projector and the resolution of the identity.}
\lectureframe{00:14:25}
\end{figure}

\begin{classroomqa}
\lecturetimestamp{00:14:38}
\audiencequestion{Should the quantity on the left be \(\overline{M}\), the average of \(M\)?}
\lecturerresponse{Yes. The trace is not the operator \(M\) itself; it is the scalar expectation value \(\overline{M}\).}
\end{classroomqa}

For a general density matrix, choose its eigenbasis:
\begin{equation}
\rho\ket{j}=\lambda_j\ket{j}.
\end{equation}
Insert a complete set of states between \(\rho\) and \(M\):
\begin{align}
\overline{M}
&=\sum_i\bra{i}\rho M\ket{i} \\
&=\sum_{i,j}\bra{i}\rho\ket{j}\bra{j}M\ket{i} \\
&=\sum_{i,j}\lambda_j\delta_{ij}\bra{j}M\ket{i} \\
&=\sum_j\lambda_j\bra{j}M\ket{j}.
\end{align}
This is an ordinary weighted average: \(\lambda_j\) is the probability of the \(j\)-th eigenstate, and \(\bra{j}M\ket{j}\) is the expectation of \(M\) in that state.

\begin{classroomqa}
\lecturetimestamp{00:18:12}
\audiencequestion{Why do both labels in the last matrix element become \(j\), and which index is still being summed?}
\lecturerresponse{Because \(\rho\) is diagonal in this basis, \(\bra{i}\rho\ket{j}=\lambda_j\delta_{ij}\), so only terms with \(i=j\) survive. The original double sum collapses to one sum, which may be called a sum over \(j\). The name of a dummy summation index has no physical significance.}
\end{classroomqa}

\section{Entropy as Missing Information}

For a classical probability distribution \(\{p_i\}\), define
\begin{equation}
S=-\sum_i p_i\log p_i,
\end{equation}
with \(0\log 0\) understood as its limiting value \(0\). If one probability equals \(1\) and all others vanish, then \(S=0\): the state is known. If all \(n\) possibilities are equally likely,
\begin{equation}
p_i=\frac{1}{n},
\qquad
S=-n\left(\frac{1}{n}\log\frac{1}{n}\right)=\log n.
\end{equation}
Thus entropy grows as the probability distribution spreads over more significant possibilities. Roughly, \(e^S\) is the effective number of comparably weighted states when natural logarithms are used.

\begin{figure}[htbp]
\centering
\includegraphics[width=0.68\textwidth]{lecture_08_figure_08.png}
\caption{Entropy for a probability distribution and the equal-probability value \(S=\log n\).}
\lectureframe{00:21:08}
\end{figure}

For a density matrix the same definition is applied to its eigenvalues:
\begin{equation}
S(\rho)=-\sum_i\lambda_i\log\lambda_i
=-\operatorname{Tr}(\rho\log\rho).
\end{equation}
The first expression is the one needed immediately; the trace form emphasizes that the answer is independent of basis.

\section{A Composite System and a Local Experiment}

Consider two subsystems, \(A\) and \(B\). A basis vector is written \(\ket{a,b}\). If \(A\) has \(N\) basis states and \(B\) has \(n\), the composite system has \(Nn\) basis states. The two systems need not be alike: one could be a two-state spin and the other a three-state atom.

The most general pure state of the pair is
\begin{equation}
\ket{\Psi}
=\sum_{a,b}\psi(a,b)\ket{a,b},
\qquad
\sum_{a,b}|\psi(a,b)|^2=1.
\end{equation}
Now let Alice measure an observable \(M_A\) that acts only on subsystem \(A\). On the composite Hilbert space the operator is \(M_A\otimes I_B\). Its matrix elements are
\begin{equation}
\bra{a',b'}M_A\otimes I_B\ket{a,b}
=\bra{a'}M_A\ket{a}\,\delta_{b'b}.
\end{equation}
The \(B\) label is passive.

Write the expectation value in full:
\begin{align}
\expect{M_A}
&=\bra{\Psi}M_A\otimes I_B\ket{\Psi} \\
&=\sum_{a',b'}\sum_{a,b}
\psi^*(a',b')\psi(a,b)
\bra{a'}M_A\ket{a}\delta_{b'b} \\
&=\sum_{a',a,b}
\psi^*(a',b)\psi(a,b)
\bra{a'}M_A\ket{a}.
\end{align}

\begin{classroomqa}
\lecturetimestamp{00:29:12}
\audiencequestion{Why are the labels on the bra primed while those on the ket are not?}
\lecturerresponse{The bra and ket are independent sums before the matrix element is evaluated, so they require separate dummy labels. Primes distinguish the bra indices from the ket indices. The operator and the inner products then determine which of those labels must coincide.}
\end{classroomqa}

Hold \(a\) and \(a'\) fixed and perform the sum over \(b\). The result is a matrix on subsystem \(A\):
\begin{equation}
\boxed{
\rho_A(a,a')
=\sum_b\psi(a,b)\psi^*(a',b)
}.
\end{equation}
This is the reduced density matrix of \(A\), or equivalently the partial trace of the pure-state projector over \(B\):
\begin{equation}
\rho_A=\operatorname{Tr}_B\ket{\Psi}\bra{\Psi}.
\end{equation}
The local expectation value becomes
\begin{equation}
\expect{M_A}
=\sum_{a,a'}\rho_A(a,a')\bra{a'}M_A\ket{a}
=\operatorname{Tr}_A(\rho_A M_A).
\end{equation}

\begin{figure}[htbp]
\centering
\includegraphics[width=0.76\textwidth]{lecture_08_figure_09.png}
\caption{The \(B\) index is summed away, leaving a matrix in the two \(A\) indices.}
\lectureframe{00:33:50}
\end{figure}

\begin{classroomqa}
\lecturetimestamp{00:35:39}
\audiencequestion{Is the local observable simply the identity on subsystem \(B\)?}
\lecturerresponse{Exactly. \(M_A\otimes I_B\) may change the \(A\) coordinate, but it neither changes nor otherwise depends on \(b\). That is why its \(B\)-space matrix element is \(\delta_{b'b}\).}
\end{classroomqa}

Every prediction for experiments confined to \(A\) is therefore contained in \(\rho_A\). Bob need not measure his subsystem. Merely declining to use the \(B\) degrees of freedom forces Alice's description into density-matrix form.

\section{A Pure Whole Can Have Mixed Parts}

The combined system is in the pure state \(\ket{\Psi}\), yet \(\rho_A\) will generally have several nonzero eigenvalues. Complete knowledge of the whole does not imply a state vector for each part. The cat and the gun may have one entangled state vector; the cat alone is then described by a reduced density matrix.

\begin{classroomqa}
\lecturetimestamp{00:38:25}
\audiencequestion{Can \(\operatorname{Tr}(M\rho)\) be replaced by \((\operatorname{Tr}M)(\operatorname{Tr}\rho)\)?}
\lecturerresponse{No. The trace of a product is generally not the product of the traces. Cyclicity permits \(\operatorname{Tr}(M\rho)=\operatorname{Tr}(\rho M)\), but it does not factorize the trace.}
\end{classroomqa}

\begin{classroomqa}
\lecturetimestamp{00:38:49}
\audiencequestion{Did we begin with only the basis state \(\ket{a,b}\), or with the most general pure state?}
\lecturerresponse{With the most general pure state of the composite system: \(\ket{\Psi}=\sum_{a,b}\psi(a,b)\ket{a,b}\). A single \(\ket{a,b}\) is only one basis vector. The surprise is that even this completely specified composite vector can give a mixed reduced state.}
\end{classroomqa}

The contrast with classical mechanics is sharp. Complete classical information about a composite system fixes each of its coordinates and hence each subsystem. A classical probability distribution may certainly be marginalized, but if the original distribution is concentrated at one phase-space point, every marginal is also concentrated at one definite value. Quantum pure states need not have that property.

\begin{classroomqa}
\lecturetimestamp{00:40:46}
\audiencequestion{How is summing over \(B\) different from integrating a classical position--velocity distribution over velocity?}
\lecturerresponse{The operations look similar, but their pure-state limits differ. A completely known classical phase-space point fixes both position and velocity, so either marginal remains definite. A completely known entangled quantum state can give a nontrivial mixed state for either subsystem. The partial trace therefore reveals a quantum failure of the classical implication from a pure whole to pure parts.}
\end{classroomqa}

\section{Product States: the Exceptional Case}

When does \(\rho_A\) remain pure? The answer is that the wave function must factorize:
\begin{equation}
\psi(a,b)=\phi(a)\chi(b).
\end{equation}
Assume \(\phi\) and \(\chi\) are separately normalized. Substitution into the reduced density matrix gives
\begin{align}
\rho_A(a,a')
&=\sum_b\phi(a)\chi(b)\phi^*(a')\chi^*(b) \\
&=\phi(a)\phi^*(a')\sum_b|\chi(b)|^2 \\
&=\phi(a)\phi^*(a').
\end{align}
Thus
\begin{equation}
\rho_A=\ket{\phi}\bra{\phi}.
\end{equation}
It has one eigenvalue \(1\), all remaining eigenvalues \(0\), and zero entropy.

\begin{figure}[htbp]
\centering
\includegraphics[width=0.74\textwidth]{lecture_08_figure_10.png}
\caption{For a factorized wave function, the \(B\)-sum reduces to the normalization of \(\chi\).}
\lectureframe{00:46:35}
\end{figure}

\begin{classroomqa}
\lecturetimestamp{00:44:17}
\audiencequestion{Is \(\phi(a)\) itself the \(A\)-system state?}
\lecturerresponse{It is the wave-function component of that state. More explicitly, \(\ket{\phi}=\sum_a\phi(a)\ket{a}\), and similarly \(\ket{\chi}=\sum_b\chi(b)\ket{b}\). Factorization means \(\ket{\Psi}=\ket{\phi}\otimes\ket{\chi}\), a product state with no entanglement between \(A\) and \(B\).}
\end{classroomqa}

For every local observable,
\begin{equation}
\expect{M_A}
=\operatorname{Tr}(\rho_A M_A)
=\bra{\phi}M_A\ket{\phi}.
\end{equation}
Subsystem \(A\) behaves exactly as a pure state \(\ket{\phi}\). For a bipartite pure state, the entropy of \(\rho_A\) therefore tests product structure: zero means a product state, while nonzero entropy means entanglement. Values near zero indicate weak entanglement; values near the maximum indicate a state far from any product across this bipartition.

\section{The Singlet: a Pure Whole and Maximal Local Ignorance}

Take the two-spin singlet
\begin{equation}
\ket{\Psi^-}
=\frac{1}{\sqrt{2}}
\left(
\ket{\uparrow\downarrow}
-\ket{\downarrow\uparrow}
\right).
\end{equation}
Its nonzero wave-function components are
\begin{equation}
\psi(\uparrow,\downarrow)=\frac{1}{\sqrt{2}},
\qquad
\psi(\downarrow,\uparrow)=-\frac{1}{\sqrt{2}},
\end{equation}
while \(\psi(\uparrow,\uparrow)=\psi(\downarrow,\downarrow)=0\).

Trace over spin \(2\). The diagonal entries are
\begin{align}
\rho_1(\uparrow,\uparrow)
&=|\psi(\uparrow,\uparrow)|^2
 +|\psi(\uparrow,\downarrow)|^2
=\frac{1}{2}, \\
\rho_1(\downarrow,\downarrow)
&=|\psi(\downarrow,\uparrow)|^2
 +|\psi(\downarrow,\downarrow)|^2
=\frac{1}{2}.
\end{align}
The off-diagonal entry is
\begin{align}
\rho_1(\uparrow,\downarrow)
&=\psi(\uparrow,\uparrow)\psi^*(\downarrow,\uparrow)
 +\psi(\uparrow,\downarrow)\psi^*(\downarrow,\downarrow) \\
&=0,
\end{align}
and its conjugate also vanishes. Hence
\begin{equation}
\rho_1
=
\begin{pmatrix}
\frac12&0\\[2pt]
0&\frac12
\end{pmatrix}
=\frac12 I.
\end{equation}

\begin{figure}[htbp]
\centering
\includegraphics[width=0.76\textwidth]{lecture_08_figure_11.png}
\caption{The four entries of the reduced density matrix for one spin in the singlet.}
\lectureframe{00:55:35}
\end{figure}

\begin{classroomqa}
\lecturetimestamp{00:55:58}
\audiencequestion{What is the entropy of the single-spin reduced state?}
\lecturerresponse{Its two eigenvalues are both \(1/2\), so
\[
S_1=-2\left(\frac12\log\frac12\right)=\log 2.
\]
This is the maximum entropy available to a two-state subsystem. Equal \emph{eigenvalues}, not equal matrix entries in an arbitrary basis, characterize the maximally mixed state.}
\end{classroomqa}

The spin expectation value vanishes along every axis:
\begin{equation}
\expect{\bm{\sigma}\cdot\hat{\bm n}}
=\operatorname{Tr}\!\left(
\rho_1\,\bm{\sigma}\cdot\hat{\bm n}
\right)
=\frac12\operatorname{Tr}
\left(\bm{\sigma}\cdot\hat{\bm n}\right)
=0.
\end{equation}

\begin{classroomqa}
\lecturetimestamp{00:58:37}
\audiencequestion{Why is the trace of \(\bm{\sigma}\cdot\hat{\bm n}\) zero for every direction?}
\lecturerresponse{Each Pauli matrix is traceless. Since \(\bm{\sigma}\cdot\hat{\bm n}=n_1\sigma_1+n_2\sigma_2+n_3\sigma_3\), linearity of the trace gives zero for every unit vector \(\hat{\bm n}\). The spin is therefore equally likely to be found parallel or antiparallel to any chosen axis.}
\end{classroomqa}

The whole pair has a definite pure state and entropy zero. Either spin by itself is maximally mixed. This is not ignorance caused by a defective preparation record; it is the local state implied by entanglement.

\section{An Equal-Amplitude Product State}

Now consider
\begin{equation}
\ket{\Phi}
=\frac12\left(
\ket{\uparrow\uparrow}
+\ket{\uparrow\downarrow}
+\ket{\downarrow\uparrow}
+\ket{\downarrow\downarrow}
\right).
\end{equation}
All four amplitudes equal \(1/2\), so normalization follows from four contributions of \(1/4\). Tracing over spin \(2\) gives
\begin{equation}
\rho_1
=
\begin{pmatrix}
\frac12&\frac12\\[2pt]
\frac12&\frac12
\end{pmatrix}.
\end{equation}
Its trace is \(1\), but
\begin{equation}
\det\rho_1
=\frac14-\frac14
=0.
\end{equation}

\begin{classroomqa}
\lecturetimestamp{01:02:23}
\audiencequestion{How do trace and determinant determine the two eigenvalues here?}
\lecturerresponse{For a \(2\times2\) matrix, the sum of the eigenvalues is the trace and their product is the determinant. Thus \(\lambda_1+\lambda_2=1\) and \(\lambda_1\lambda_2=0\), which forces the pair \(\{1,0\}\). The reduced state is pure and its entropy is zero.}
\end{classroomqa}

\begin{classroomqa}
\lecturetimestamp{01:03:18}
\audiencequestion{What are the two unentangled spins actually doing?}
\lecturerresponse{They are both aligned along the positive \(x\)-axis. The answer can be recognized by factorization,
\[
\ket{\Phi}
=\frac{\ket{\uparrow}+\ket{\downarrow}}{\sqrt2}
\otimes
\frac{\ket{\uparrow}+\ket{\downarrow}}{\sqrt2},
\]
or established experimentally by computing the Pauli expectation values.}
\end{classroomqa}

For the \(z\)-component,
\begin{equation}
\sigma_3=
\begin{pmatrix}
1&0\\[2pt]
0&-1
\end{pmatrix},
\qquad
\operatorname{Tr}(\rho_1\sigma_3)=0.
\end{equation}
For the \(x\)-component,
\begin{equation}
\sigma_1=
\begin{pmatrix}
0&1\\[2pt]
1&0
\end{pmatrix},
\end{equation}
and
\begin{align}
\rho_1\sigma_1
&=
\begin{pmatrix}
\frac12&\frac12\\[2pt]
\frac12&\frac12
\end{pmatrix}
\begin{pmatrix}
0&1\\[2pt]
1&0
\end{pmatrix}
=
\begin{pmatrix}
\frac12&\frac12\\[2pt]
\frac12&\frac12
\end{pmatrix}, \\
\operatorname{Tr}(\rho_1\sigma_1)&=1.
\end{align}

\begin{figure}[htbp]
\centering
\includegraphics[width=0.70\textwidth]{lecture_08_figure_03.png}
\caption{The \(\sigma_1\) matrix calculation for the equal-amplitude reduced state.}
\lectureframe{01:05:13}
\end{figure}

Thus \(\expect{\sigma_1}=1\) while \(\expect{\sigma_2}=\expect{\sigma_3}=0\). Perturbing the amplitudes away from exact factorization generally changes the eigenvalues from \(\{1,0\}\) to two nonzero values. The resulting small entropy measures a small amount of entanglement.

\section{Whose Entropy Is It?}

The entropy just computed belongs to a subsystem. In the singlet example,
\begin{equation}
S(\rho_{12})=0,
\qquad
S(\rho_1)=S(\rho_2)=\log 2.
\end{equation}
The whole is pure while each part is mixed. Entanglement entropy therefore does not obey the naive additive rule \(S_{12}=S_1+S_2\). For a bipartite pure state the two reduced density matrices have the same nonzero eigenvalues, so their entropies are equal even though the entropy of the whole vanishes.

\begin{classroomqa}
\lecturetimestamp{01:09:03}
\audiencequestion{Is every density matrix ultimately an entanglement density matrix?}
\lecturerresponse{Operationally, one may describe ordinary uncertainty without tracking what carries the missing record. If the preparation apparatus, environment, or observer is included, that uncertainty can often be viewed as entanglement with unobserved degrees of freedom. Mathematically the idea is exact: every density matrix admits a purification on a larger Hilbert space, although that purification does not uniquely identify a physical preparation history.}
\end{classroomqa}

\section{From Discrete Classical Motion to Quantum Motion}

We now leave entanglement temporarily and ask how states change with time. For a classical system with finitely many configurations, reversible evolution is a permutation. Each state has one successor and one predecessor; the motion decomposes into cycles.

\begin{figure}[htbp]
\centering
\includegraphics[width=0.70\textwidth]{lecture_08_figure_12.png}
\caption{Classical reversible evolution represented by arrows among discrete configurations.}
\lectureframe{01:11:40}
\end{figure}

\begin{classroomqa}
\lecturetimestamp{01:11:51}
\audiencequestion{Why must the time steps also be discrete if the configurations are discrete?}
\lecturerresponse{A strictly classical two-state model contains no intermediate configurations through which a continuous path could pass. Turning a physical coin slowly introduces its orientation as an additional continuous coordinate, so it is no longer a two-configuration model. A quantum two-state system is different: its Hilbert space contains a continuous family of superpositions even though an orthonormal basis has only two members.}
\end{classroomqa}

A spin-\(\tfrac12\) state can rotate continuously from \(\ket{\uparrow}\) toward \(\ket{\downarrow}\) through superpositions. Finite dimension does not imply a finite set of state vectors. This allows continuous time evolution in a discrete quantum system.

Reversible microscopic dynamics also preserves distinctions. Two different classical initial conditions do not merge into one microscopic state, because that would erase which initial condition occurred. In Hamiltonian mechanics this invertibility is accompanied by Liouville's theorem: phase-space volume is preserved.

\begin{classroomqa}
\lecturetimestamp{01:15:26}
\audiencequestion{What about nonlinear bifurcations in which several initial conditions approach the same steady state?}
\lecturerresponse{Attractors occur in effective dissipative descriptions, where environmental degrees of freedom have been discarded. Fundamental microscopic Hamiltonian evolution remains invertible; information may be transferred to unobserved variables, but distinct complete states do not literally become one complete state.}
\end{classroomqa}

The quantum statement is unitarity. Orthogonal initial states evolve into orthogonal final states, and more generally every inner product is preserved.

\begin{classroomqa}
\lecturetimestamp{01:16:50}
\audiencequestion{Must the two states be subjected to exactly the same influence? What happens to initial \(z\)-up and \(z\)-down states in an \(x\)-directed magnetic field?}
\lecturerresponse{The same evolution operator must act on both states. In the transverse field the two spins precess, but their Bloch vectors remain opposite and the state vectors remain orthogonal. Applying different external histories would amount to comparing different evolution operators.}
\end{classroomqa}

\section{Linearity and Unitarity}

The first principle of quantum time evolution is linearity:
\begin{equation}
\ket{\psi(t)}=U(t)\ket{\psi(0)}.
\end{equation}
The operator depends on the elapsed time and satisfies
\begin{equation}
U(0)=I.
\end{equation}
Like the linearity of quantum mechanics itself, this rule is justified by its empirical success.

\begin{figure}[htbp]
\centering
\includegraphics[width=0.70\textwidth]{lecture_08_figure_13.png}
\caption{The state at time \(t\) is obtained by a linear operator \(U(t)\) acting on the initial state.}
\lectureframe{01:19:00}
\end{figure}

\begin{classroomqa}
\lecturetimestamp{01:19:01}
\audiencequestion{What must \(U(t)\) be at \(t=0\)?}
\lecturerresponse{It must be the identity. At zero elapsed time the state is unchanged for every possible initial vector.}
\end{classroomqa}

The second principle is preservation of inner products:
\begin{equation}
\braket{\phi(t)}{\psi(t)}
=\braket{\phi(0)}{\psi(0)}.
\end{equation}
When \(\phi=\psi\), this preserves normalization and hence total probability. When the initial states are orthogonal, it preserves perfect distinguishability. For arbitrary states it preserves their entire geometric relation.

The bra evolves according to
\begin{equation}
\bra{\phi(t)}
=\bra{\phi(0)}U^\dagger(t),
\end{equation}
so
\begin{align}
\braket{\phi(t)}{\psi(t)}
&=\bra{\phi(0)}U^\dagger(t)U(t)\ket{\psi(0)} \\
&=\braket{\phi(0)}{\psi(0)}.
\end{align}
Since this equality holds for every pair of initial vectors,
\begin{equation}
\boxed{U^\dagger(t)U(t)=I}.
\end{equation}

\begin{figure}[htbp]
\centering
\includegraphics[width=0.74\textwidth]{lecture_08_figure_04.png}
\caption{Ket evolution, bra evolution, and the inner-product condition leading to \(U^\dagger U=I\).}
\lectureframe{01:23:30}
\end{figure}

\begin{classroomqa}
\lecturetimestamp{01:22:10}
\audiencequestion{Why does the bra acquire \(U^\dagger\) rather than \(U\)?}
\lecturerresponse{Taking the Hermitian conjugate reverses the order of factors and replaces each operator by its Hermitian conjugate:
\[
\bigl(U\ket{\phi}\bigr)^\dagger
=\bra{\phi}U^\dagger.
\]
That is the dual action required for consistent inner products.}
\end{classroomqa}

\begin{classroomqa}
\lecturetimestamp{01:24:34}
\audiencequestion{If \(U\) were just a complex number, what would \(U^\dagger U=1\) mean?}
\lecturerresponse{It would mean \(|U|=1\), so \(U=e^{i\theta}\) lies on the unit circle. A unitary matrix is the operator analogue of a pure phase: it preserves lengths and inner products.}
\end{classroomqa}

\section{Infinitesimal Evolution and the Hamiltonian}

Let \(\epsilon\) be a very small time interval. Since \(U(0)=I\), write the first-order expansion
\begin{equation}
U(\epsilon)=I-i\epsilon h+O(\epsilon^2).
\end{equation}
The minus sign and the factor of \(i\) are convenient conventions. Taking the Hermitian conjugate gives
\begin{equation}
U^\dagger(\epsilon)
=I+i\epsilon h^\dagger+O(\epsilon^2).
\end{equation}
Unitarity then requires
\begin{align}
I
&=U^\dagger(\epsilon)U(\epsilon) \\
&=\left(I+i\epsilon h^\dagger\right)
\left(I-i\epsilon h\right)+O(\epsilon^2) \\
&=I+i\epsilon(h^\dagger-h)+O(\epsilon^2).
\end{align}
The first-order term must vanish:
\begin{equation}
\boxed{h^\dagger=h}.
\end{equation}
The infinitesimal generator is Hermitian.

\begin{figure}[htbp]
\centering
\includegraphics[width=0.76\textwidth]{lecture_08_figure_14.png}
\caption{First-order unitarity forces the infinitesimal generator to be Hermitian.}
\lectureframe{01:29:35}
\end{figure}

In the temporary convention above, \(h\) has units of angular frequency. Restoring physical units, define the energy operator
\begin{equation}
H=\hbar h.
\end{equation}
This Hermitian operator is the Hamiltonian, and its eigenvalues are energy levels.

Apply the infinitesimal operator to a state:
\begin{equation}
\ket{\psi(\epsilon)}
=\left(I-\frac{i\epsilon}{\hbar}H\right)\ket{\psi(0)}
+O(\epsilon^2).
\end{equation}
Subtract the initial state, divide by \(\epsilon\), and take the limit:
\begin{equation}
\frac{\mathrm d}{\mathrm dt}\ket{\psi(t)}
=-\frac{i}{\hbar}H\ket{\psi(t)}.
\end{equation}
Equivalently,
\begin{equation}
\boxed{
i\hbar\frac{\mathrm d}{\mathrm dt}\ket{\psi(t)}
=H\ket{\psi(t)}
}.
\end{equation}
This is the general Schr\"odinger equation.

\begin{figure}[htbp]
\centering
\includegraphics[width=0.76\textwidth]{lecture_08_figure_15.png}
\caption{The infinitesimal difference quotient becomes the time derivative of the state.}
\lectureframe{01:32:18}
\end{figure}

The Hamiltonian depends on the physical setup. It may change with time if an external magnetic field or another control changes. When the setup is time independent, \(H\) is a fixed Hermitian operator.

\begin{classroomqa}
\lecturetimestamp{01:34:30}
\audiencequestion{Does preservation of inner products imply that initially unentangled systems must remain unentangled?}
\lecturerresponse{No. Unitarity preserves inner products between complete state vectors, not the product structure of a chosen partition. An interaction Hamiltonian can carry a product state into an entangled state. If unitaries could not do this, measurements could not correlate a system with an apparatus.}
\end{classroomqa}

\section{The Units of Entropy}

\begin{classroomqa}
\lecturetimestamp{01:35:08}
\audiencequestion{Does entanglement have a physical unit such as joules?}
\lecturerresponse{Entanglement entropy is an information measure, not an energy. With natural logarithms it is measured in nats; with base-two logarithms it is measured in bits. Both are dimensionless counting conventions.}
\end{classroomqa}

Thermodynamic notation can obscure this point because the Boltzmann constant converts between temperature measured in kelvins and energy:
\begin{equation}
\expect{K}_{\mathrm{monatomic}}
=\frac32 k_{\mathrm B}T.
\end{equation}
The reversible heat relation is
\begin{equation}
\delta Q_{\mathrm{rev}}=T\,\mathrm dS_{\mathrm{th}}.
\end{equation}
If the information entropy is
\begin{equation}
S_{\mathrm{info}}=-\sum_i p_i\log p_i,
\end{equation}
then the conventional thermodynamic entropy is
\begin{equation}
S_{\mathrm{th}}=k_{\mathrm B}S_{\mathrm{info}}.
\end{equation}
Setting \(k_{\mathrm B}=1\) measures temperature in energy units and leaves entropy dimensionless. Boltzmann's constant is then recognized as a conversion factor rather than a new kind of entanglement unit. The historical aside connecting Boltzmann's statistical ideas with Einstein's 1905 work on Brownian motion does not alter the information-theoretic definition.

\section{Energy Eigenstates and Phase}

\begin{classroomqa}
\lecturetimestamp{01:39:54}
\audiencequestion{Why is the Hermitian generator \(H\) identified with energy?}
\lecturerresponse{Unitarity alone proves only that the generator is Hermitian. Its identification with energy is an additional physical statement: energy generates time translations. For a time-independent Hamiltonian, the same operator that governs evolution is conserved, and its eigenvalues are the stationary energy levels. Concrete Hamiltonians supply the experimental content of that identification.}
\end{classroomqa}

The relation between Planck's constants and frequency is
\begin{equation}
E=h\nu=\hbar\omega,
\qquad
\omega=2\pi\nu,
\qquad
h=2\pi\hbar.
\end{equation}

\begin{classroomqa}
\lecturetimestamp{01:40:30}
\audiencequestion{Should the evolution equation contain \(h\), \(\hbar\), or an extra factor of \(2\pi\)?}
\lecturerresponse{Using angular frequency \(\omega\), the equation is \(E=\hbar\omega\), and the Schr\"odinger equation contains \(\hbar\). Using ordinary frequency \(\nu\), the equivalent relation is \(E=h\nu\). The factor \(2\pi\) is entirely accounted for by \(\omega=2\pi\nu\) and \(h=2\pi\hbar\).}
\end{classroomqa}

Suppose the initial state is an energy eigenvector:
\begin{equation}
H\ket{\psi(0)}=E\ket{\psi(0)}.
\end{equation}

\begin{classroomqa}
\lecturetimestamp{01:43:18}
\audiencequestion{If the state remains an energy eigenvector, does that mean it does not change?}
\lecturerresponse{It may acquire a time-dependent complex factor and remain an eigenvector with the same eigenvalue. Unitarity restricts that factor to a phase. Thus the ray is stationary even while the ket rotates in phase.}
\end{classroomqa}

Write
\begin{equation}
\ket{\psi(t)}=f(t)\ket{\psi(0)}.
\end{equation}

\begin{figure}[H]
\centering
\includegraphics[width=0.70\textwidth]{lecture_08_figure_05.png}
\caption{An energy eigenstate remains along the same Hilbert-space direction up to a scalar factor \(f(t)\).}
\lectureframe{01:43:54}
\end{figure}

Substitution into the Schr\"odinger equation gives
\begin{align}
i\hbar\frac{\mathrm d f}{\mathrm dt}\ket{\psi(0)}
&=f(t)H\ket{\psi(0)} \\
&=E f(t)\ket{\psi(0)}.
\end{align}
Cancel the common vector:
\begin{equation}
\frac{\mathrm d f}{\mathrm dt}
=-\frac{iE}{\hbar}f(t).
\end{equation}
With \(f(0)=1\),
\begin{equation}
f(t)=e^{-iEt/\hbar},
\qquad
\ket{\psi(t)}=e^{-iEt/\hbar}\ket{\psi(0)}.
\end{equation}

\begin{figure}[H]
\centering
\includegraphics[width=0.76\textwidth]{lecture_08_figure_16.png}
\caption{The energy eigenvalue equation and the resulting phase evolution.}
\lectureframe{01:46:10}
\end{figure}

\begin{classroomqa}
\lecturetimestamp{01:45:39}
\audiencequestion{What is the coefficient of \(t\) in the phase \(e^{-iEt/\hbar}\)?}
\lecturerresponse{It is the angular frequency,
\[
\omega=\frac{E}{\hbar}.
\]
Equivalently \(E=\hbar\omega\). A real exponential coefficient would describe a growth or decay rate; the imaginary exponent here describes unitary phase rotation.}
\end{classroomqa}

This phase-frequency relation and the Schr\"odinger equation are two forms of the same law of time evolution. The equation is general; a free particle, a particle in a potential, and a spin in a magnetic field differ through their Hamiltonians. Choosing and solving those Hamiltonians is the next step.
